\chapter{第 15 章分层答案}
\label{app:ch15-solutions}

\section*{口述概念}
\begin{enumerate}[leftmargin=*]
  \item seed 只使同一伪随机序列可重放。程序可以稳定地重放错误分布、错误支付函数或错误单位；正确性还需要解析或独立计算参照、误差率与失败测试。
  \item 模型误差来自目标分布不对；离散化偏差来自时间步或网格；抽样误差来自有限随机样本；数值误差来自浮点和病态计算；实现错误来自代码没有计算定义的对象。增加 $N$ 主要降低第三项。
  \item Monte Carlo 在给定模型 $P$ 下近似 $\mathbb E_P[g(X)]$。IID bootstrap 条件于观测样本，从经验分布 $\hat F_n$ 重抽样，近似统计量的条件抽样分布；它要求 IID 样本或另有理论保证的等价设计，并要求统计量满足相应正则性。交换性本身不保证独立。
\end{enumerate}

\section*{笔试推导}
\begin{enumerate}[leftmargin=*]
  \item
  \[
    \operatorname{Var}(N^{-1}\textstyle\sum_iY_i)
    =N^{-2}\sum_{i,j}\operatorname{Cov}(Y_i,Y_j).
  \]
  IID 时非对角项为零，留下 $N\sigma^2/N^2=\sigma^2/N$。相关样本的非对角项保留，标准误不再简单等于 $\sigma/\sqrt N$。
  \item 展开
  $V(\beta)=\operatorname{Var}(g)-2\beta\operatorname{Cov}(g,h)+\beta^2\operatorname{Var}(h)$，令
  $V'(\beta)=0$ 得 $\beta^*=\operatorname{Cov}(g,h)/\operatorname{Var}(h)$。最小方差为
  $\operatorname{Var}(g)(1-\rho^2)$。对 $g(U)=U^2,h(U)=U$，解析最优系数为 $1$，调整后方差 $1/180$，理论方差缩减倍数为 $16$。
  \item 三次抽样和的频数依次为
  $(3,4,5,6,7,8,9,10,12)\mapsto(1,3,3,4,6,3,3,3,1)$。除以 3 得到均值分布；按频数计算一阶、二阶中心矩，得到 $\mathbb E_*(\bar X^*)=7/3$、$\operatorname{Var}_*(\bar X^*)=14/27$。
\end{enumerate}

\section*{数值编程}
\begin{enumerate}[leftmargin=*]
  \item 三种目标均为 $1/3$，单次贡献方差分别为 $4/45,1/180,1/72$。
  \item 从 $0.5$ 改成 $0.6$ 时均值增加 $0.1$，方差不变。
  \item 共 27 个，条件均值 $7/3$，条件方差 $14/27$；增加重抽样次数不增加原始样本信息。
\end{enumerate}

\section*{研究判断}
\begin{enumerate}[leftmargin=*]
  \item 逐日 IID 重抽样破坏自相关、波动聚集和路径依赖，通常低估 Sharpe 或回撤不确定性。应说明 block 长度依据，比较移动块/平稳 bootstrap、时间外样本和模拟覆盖率；结构变化明显时降低区间解释强度。
  \item 重要性权重巨大说明 $Q$ 与目标重要区域匹配不佳，估计可能由极少路径支配。应检查支撑覆盖、权重尾部、最大权重、有效样本量、二阶矩是否有限，并比较多个提议分布与独立 seeds。
  \item 只对最终因子 bootstrap 忽略了从五百个候选中选择这一随机步骤，区间是选择后条件区间而非完整研究流程的不确定性。需要在每次重抽样中重做筛选，或使用嵌套/样本分割方案，并把多重检验与样本外验证纳入报告。
\end{enumerate}

\section*{基础衔接：独立计算}
贡献 $x^2/q(x)=1/3$ 几乎处处恒定，方差为零。这里提议密度的归一化常数已经用到了目标积分；一般未知积分不能凭此免费构造可抽样的零方差密度。
